\section{Recomendaciones}

\begin{enumerate}
	\item Instalar un sensor de presión diferencial entre el laboratorio y la zona adyacente, con un set-point de control entre +12.5 Pa y +25 Pa, para garantizar y monitorear la presurización positiva de forma continua.
	\item Filtrar el aire de impulsión con filtros de mínimo eficiencia MERV 13-14. Para aplicaciones con riesgo biológico elevado, considerar filtros HEPA H13 o H14 según la normativa sanitaria local.
	\item Sellar adecuadamente puertas, ventanas, pasos de instalaciones y cualquier abertura no controlada del recinto, con burletes y umbrales herméticos, para reducir las fugas no controladas y facilitar el mantenimiento de la presión positiva.
	\item Ubicar las tres rejillas de exfiltración en paredes opuestas o perpendiculares a la impulsión, de modo que se favorezca un flujo cruzado que minimice zonas de estancamiento.
	\item Ejecutar un modelo CFD con las condiciones de contorno entregadas para validar la distribución de velocidades, la edad del aire y el campo de presiones en el recinto, antes de la construcción final.
	\item Prever un margen de sobredimensionamiento del 10 \% al 20 \% entre el caudal de impulsión y cualquier caudal de extracción localizada (campanas, puntos de captación), de modo que la presión positiva general del recinto no se vea comprometida.
\end{enumerate}
